\documentclass[border=3pt]{standalone}

% Tikz
\usepackage{tikz}
\usetikzlibrary{calc}

% Sizes
\def\cellheight{30pt}
\def\fieldlinewidth{1pt}
\pgfmathsetlengthmacro{\scalefactor}{5pt}

% Tikz Set
\tikzset{
    instr_field/.style={
        draw,
        rectangle,
        line width = \fieldlinewidth,
        minimum height = \cellheight,
        rounded corners=2pt,
        fill=teal!15,
        text=blue!80!black,
    },
    byte_values/.style={below, font=\tiny, text=red!80!black},
    bytes/.style={above, font=\scriptsize, text=purple!90!black},
    every text node part/.style={align=center}
}

% Instruction Fields (renamed for Computer Science semantics, values changed)
\def\fields{
    flag/31/29/101,
    opcode/28/25/1100,
    type/24/24/0,
    func/23/20/0111,
    ext/19/19/1,
    src1/18/15/0010,
    src2/14/11/1111,
    shift/10/8/101,
    imm5/7/3/01010
}

\begin{document}
    \begin{tikzpicture}
        % Widths
        \foreach \field/\high/\low in \fields {
            \pgfmathsetmacro{\width}{(\high-\low + 1) * \scalefactor}
            \expandafter\xdef\csname \field w\endcsname{\width mm}
        }

        % Instruction Fields
        \foreach \field/\high/\low/\value [count=\i from 0] in \fields {
            \ifnum\i=0
                \node[
                    instr_field,
                    anchor=west,
                    minimum width=\csname\field w\endcsname
                ] (tmp\i) {\field};
            \else
                \pgfmathsetmacro{\prev}{int(\i-1)}
                \node[
                    instr_field,
                    anchor=west,
                    minimum width=\csname\field w\endcsname
                ] (tmp\i) at ($(tmp\prev.east)-(\fieldlinewidth,0)$) {\field};
            \fi

            \ifnum\high=\low
                \node[bytes] at (tmp\i.north) {\high};
            \else
                \node[bytes] at (tmp\i.north) {\high:\low};
            \fi

            \node[byte_values] at (tmp\i.south) {\value};
        }

        % Example instruction name
        \node[anchor=east, font=\small\bfseries, text=teal!80!black, xshift=-32pt] at (tmp0) {LOGIC-OP};
    \end{tikzpicture}
\end{document}